\section{Git}
Do kontroli wersji aplikacji webowej i aplikacji na STM32 został wybrany Git \cite{Git}. Git jest systemem kontroli wersji, który umożliwia zarządzanie historią zmian w projekcie. Jest niezwykle przydatny w celu podejrzenia wszelkich zmian pomiędzy commit'ami.
\subsection{Podstawowe cechy Git}
Git działa w oparciu o repozytorium, w którym przechowywane są wszystkie wersje plików. Repozytorium może być zarówno lokalne jak i online, w tym przypadku wybrano repozytorium online na portalu GitHub, ponieważ ułatwiło to pracę w zespole jak i na wielu urządzeniach. Główne narzędzia dostępne w Git, które ułatwiają pracę:
\begin{enumerate}
    \item \textbf{Commit:} zapisanie zmian plików do historii. Commit powinien również zawierać opis, który obrazuje dokonane zmiany.
    \item \textbf{Branch (gałąź):} wersja projektu rozwijana równolegle z innymi. Dzięki gałęziom możemy rozwijać różne funkcje równocześnie bez zakłócania pracy nad główną wersją projektu. Jest to dobra praktyka w trakcie rozwijania nowej funkcjonalności, ponieważ nie dodajemy niestabilnych elementów na główny branch. Nazwą głównego brancha na portalu GitHub jest \textit{main}.
    \item \textbf{Merge:} połączenie zmian z jednej gałęzi z inną. Pozwala na przeniesienie efektów pracy z jednej gałęzi do głównej lub między innymi gałęziami. Było to przydatne w celu dodania 
 skończonych i sprawdzonych nowych funkcjonalności na główny branch.
    \item \textbf{Tags (tagi):} to etykiety przypisywane do danej wersji projektu. Tagi są przydatne do szybkiego nawigowania między wersjami projektu. Tagi zostały przede wszystkim wykorzystane do tworzenia stabilnych wersji, które przyczyniły się do testowania czy nowe oprogramowanie wpłynęło na negatywne działanie całego projektu.
\end{enumerate}
\subsection{Praca sieciowa}
Git przechowuje całą historię zmian w plikach w postaci serii commit'ów i przechowuje zmiany w postaci zestawu różnic (ang. \textit{diff}) między kolejnymi wersjami plików. Cała historia znajduje się na serwerze, którą można pobrać na lokalny komputer i pracować bez połączenia z siecią. Połączenie z siecią jest potrzebne do pobrania najnowszej wersji projektu \textit{git pull} lub \textit{git fetch} i do wysyłania własnych zmian \textit{git push}.
\subsection{Podsumowanie}
Git to przydatne narzędzie, które wprowadza porządek do projektu. Ułatwia kontrole nad wersjami aplikacji i umożliwia równoległą pracę nad wieloma funkcjonalnościami oraz zapewnia bezpieczeństwo dzięki historii zmian. W przypadku problematycznych nowych funkcjonalności można się cofnąć do poprzedniej stabilnej wersji.

